\documentclass[11pt]{article}
\usepackage{amsmath,amssymb,hyperref}
\title{A route probe can consume the capacity needed to deliver both reports}
\author{}\date{}
\begin{document}\maketitle
This authored civil inspection fixture deliberately couples route learning
to an outstanding radio service duty. After two inspections and durable
local writes (ticks 0--2), report A resides on relay slots 0 and 1, and
report B on slots 1 and 2. Neither report has arrived at base. A later
permanent silent relay does not erase the other actual durable copies.
Each live slot can forward at most one report in the single forwarding
round. An assignment succeeds only when both real copies reach base.
The remaining worker starts at junction J at tick two and must actually
arrive at base. Its two egress ducts are one-way once entered. Either
distal exit can close, and J does not see that closure. An intact
advisory link can report the closure and live relay set; the compound
case loses that link. Pure relay status or completed report service
does not by itself establish egress.

The next worker can travel one tick from J to probe P, or two ticks
along another branch to probe Q. A completed probe costs one tick;
both return which exit is blocked and which of the two possible
outer relays failed on the admitted stable source. P uses the same
power feed as slot 1 and disables that slot until the sortie ends.
Q has an independent feed and leaves the slot available. These
are stipulated source mechanisms, not inferences from underground
radio behavior. After the probe, the worker returns to J and can
take a viable duct in two ticks. A forwarding decision itself takes
one full tick and is issued from J or a reached probe station; it is
not a move. Local probes are not received remotely. The tester acts
from current position, local adjacency, observed evidence, and its
own previous actions; the actual blocked exit, loss and custody
remain evaluator state until a permitted read or effect exposes them.
Synchronized mission ticks, reliable acquired evidence, and explicit
source changes are part of the fixture.

\paragraph{One common physical history.}
Taking each local handler at its own advertised contract gives a
false positive at eight ticks: P can identify the route, and the
relay handler, given the reported surviving state but not the
feed diversion, can schedule both reports. Together they cannot.
For example, when slot 0 failed, A's only remaining copy is on
slot 1, which P actually disconnected. When slot 2 failed the
symmetric failure is B's. Without inspecting, no one assignment
survives both statuses: A needs slot 1 under failure 0, and B needs
slot 1 under failure 2. Full labels cannot restore the unavailable slot.
Using Q permits the two assignments (1,2) and (0,1). An independent
feed at P supplies the same capability two ticks sooner, but this
spends a physically different resource.

In the compound stable case Q permits both reports by tick six and
egress by ten. No plan can deliver both by five in that fixture,
even allowing a longer egress deadline. No probe or P cannot
satisfy both duties, so any successful plan must reach Q (two
ticks), read for one, and forward for one, all after tick two;
the return and egress take another four. If only the distal
closure occurs with no late radio loss, forwarding first and then
probing the route avoids the P feed conflict. The same is true
whenever the initial observation already allows a common forward
assignment and P's effect is subsequent to actual settlement.
With Q's branch closed, or with a second physical silent relay
after Q's read that removes the still-needed slot 1, the scoped
case refuses rather than leaving a report undelivered.

Persistence of the feed diversion is material. If leaving P
physically reconnects slot 1 before another forward decision,
the controller reads at P, returns to J, then transmits. This
alternative source law reaches base by eight and delivers by
six. The earlier impossibility is not a claim about dispensable
or freely reversible probes.

\paragraph{Exact information requirement of a single read.}
At the read cut, separate eligibility by actual same-history effect.
For each observation cohort, a feasible assignment uses a distinct
slot for each report, surviving in all worlds of that cohort under
the reservation incurred by making this observation. A common
egress path must also exist. For the two paired scenarios, two
distinct observations suffice. If distal closure and the failed
outer relay vary independently, there are four distinguishable
classes: merging any two with different closures loses a common
egress route, and merging two with the same closure but different
failed relays loses a common assignment. The two families have
the same individual support for the closure and the loss variables.
A proof relying on one field as a proxy for the other cannot
silently transfer to the four-scenario family.

\paragraph{An observed access change selects a different tester.}
A separate four-world control allows the far probe's drive to be
severed. In those worlds only, the near probe is independently powered.
Neither the egress failure nor the failed outer relay is determined
by the drive status. Approaching the first waypoint of Q reveals
whether the onward segment actually exists. When it does, proceed
to Q; when it does not, return and use P. With this lawful local
branch, exact planning and a direct written-down policy both deliver
by seven and reach base by ten. Neither selection of one station
before seeing access can serve all four worlds. Choosing P
immediately is unsafe for timely service in the shared-feed worlds;
choosing Q immediately encounters the physical obstruction in the
other worlds. If the access observation itself additionally determines
the exit and relay loss, a smaller two-world control dispenses with
the powered acquisition and completes by seven. That stronger
correlation is not available in the four-world control.

\paragraph{Comparison and scope.}
Nominal, closure-only with intact advisory, and loss-of-advisory-only
cases finish by five. The exact local-history tree runs actual
effect traces and rejects uncompleted reports and trapped arrivals.
The matched factored planner certifies common assignment and egress
per acquisition cohort under the same source effects, testing both
send-before-read and read-before-send. It ties on all 256
instantiated complete comparisons. The factored benchmark covers
common safe access to a chosen probe inside the initial observation
cell and does not claim completeness for move-induced discovery
of access on the route to that probe. Three Lean results check
source action feasibility and physical shared-feed failure.
The timing and route checks do not alter the occupation criteria
of the earlier two-hop corridor result.

The closest operational starting points include the
\href{https://superintelligence.jpl.nasa.gov/public/documents/papers/PlanRob-2020-Planning-and-Optimization-for-Multi-Robot-Planetary-Cave-Exploration.pdf}{JPL cave return-and-rendezvous planning example}
and ordinary matching on physically eligible relay slots. This
is an authored calibration with a matched factored reference, not
a reimplementation of their general system, an advantage over a
matched robust resource-constrained planner, a true sampled
underground law, a frozen independent transfer, or acceptance of I2.
The input technology that would substantiate a stable last read and
safe simultaneous forwarding slots is still an explicit obligation
for any independent deployment.
\end{document}
